\section{Data Preprocessing and Exploration}\label{sec:data}

We use only the four official files in \texttt{ProblemC\_Music\_data/}.
No external data are introduced.

\subsection{Files and roles}
\begin{itemize}
  \item \texttt{influence\_data.csv}: directed influencer$\to$follower edges, with \texttt{main\_genre} and \texttt{active\_start} (used later for networks; here mainly as the genre dictionary).
  \item \texttt{data\_by\_artist.csv}: artist-level mean Spotify features (primary input to Model~II).
  \item \texttt{full\_music\_data.csv}: song-level features ($\approx 98{,}340$ tracks; reserved for finer temporal analyses).
  \item \texttt{data\_by\_year.csv}: yearly mean features from 1921--2020 (reserved for evolution plots).
\end{itemize}

\subsection{Cleaning choices for Model~II}
\begin{enumerate}
  \item Join \texttt{data\_by\_artist} to genre labels via \texttt{artist\_id}, where genres are taken from \texttt{influence\_data} (we do \emph{not} redefine Pop/Rock or other genres).
  \item Drop artists with missing genre or genre \texttt{Unknown}.
  \item Select $p=11$ numeric features:
  danceability, energy, valence, tempo, loudness, acousticness, instrumentalness, liveness, speechiness, duration\_ms, popularity.
  \item Apply column-wise $z$-score standardization before computing Euclidean distances (see Model~II).
\end{enumerate}
After cleaning we retain $n=5599$ artists with both features and a genre label, covering 19 genres (Pop/Rock is the largest class).
These steps are reproduced by \texttt{code/preprocess\_data.py}, which writes \texttt{artists\_features\_with\_genre.csv} and a short summary JSON for auditability.
